% =====================================================================
%  tr01c —— §3 讲义（原文 pp.12--13）
%  归属：tr-01。kappa 类、psi 类、定理 3（p.12）、"Volume results overview"、
%  定理 4 与定理 5（p.13）；图 5（p.12）。本块无编号公式。
% =====================================================================
% [原文 p.12]
以及沿带节点 Riemann 曲面的相对对偶化层。各个纤维上的双曲度量为
$(\mathrm{Ker}\,\dd\pi)$ 在 $\overline{\cC}$ 上提供一个线丛度量；它虽然不光滑，
但对 Chern 形式 $\mathbf{c}_1$ 的计算而言已足够正则。由沿纤维积分给出的 kappa
形式／上同调类
\[
\kappa_k=\int_{\pi^{-1}(\{R\})}\mathbf{c}_1^{k+1}
\]
是模空间几何的基础。

kappa 类的几何与代数在 Carel Faber 的讲义中讨论。

对 Chern 形式的显式计算与积分给出下面的结果。

% [Tutor Note: p.12 的 $\kappa_k=\int_{\pi^{-1}(\{R\})}\mathbf{c}_1^{k+1}$ 必须
%  读成"沿纤维的积分"（纤维积分／pushforward），而不是普通积分：纤维是复曲线，
%  普通积分只能容下次数 $\le 2$ 的形式，而 $\mathbf{c}_1^{k+1}$ 次数是
%  $2k+2$，只有把沿纤维方向的形式积掉、留下底空间上的 $2k$-形式才讲得通。故
%  $k=1$ 时 $\kappa_1=\pi_*(\mathbf{c}_1^2)$ 是 $\M$ 上的 $2$-形式。原文这里的
%  $\mathbf{c}_1$ 是粗体（已裁 1200 dpi 复核），表示 Chern 形式而非类。]

% [原文 p.12] 图 5
\begin{figure}[htbp]
\centering
\includegraphics[width=0.55\textwidth]{fig05.png}
\figtitle{5}{万有曲线 $\overline{\cC}/\MCG_{\overline{\cC}}$ 的纤维上的切向量。}
\end{figure}

\begin{thm}{3} (W)，[Wlp90]。对纤维上的双曲度量，$2\pi^2\kappa_1=\omega$ 在
$\M$ 上逐点成立，并在 $\Mbar$ 上于上同调意义下成立。\end{thm}

% [Tutor Note: p.12 定理 3 的常数 $2\pi^2$ 是全书最容易抄错的地方。用
%  $V_{1,1}(L)=\pi^2/6+L^2/24$ 反向核对最省事：$(g,n)=(1,1)$ 时
%  $3g-3+n=1$，定理 4 给 $C_{(1)}=\int_{\Mbar_{1,1}}\psi_1=1/24$（分子
%  $2^{\delta_{1g}\delta_{1n}}$ 与分母 $2^{|\alpha|}\alpha!$ 当 $|\alpha|=1$ 时
%  恰好约成 $1$），而 $C_{(0)}=2\int_{\Mbar_{1,1}}\omega=\pi^2/6$，即
%  $\int_{\Mbar_{1,1}}\omega=\pi^2/12=2\pi^2\cdot\frac{1}{24}$，取的正是
%  $\int_{\Mbar_{1,1}}\kappa_1=1/24$ 与 $\omega=2\pi^2\kappa_1$。若把 $2\pi^2$
%  误抄成 $\pi^2$，常数项立刻变成 $\pi^2/12$，与已知的 $\pi^2/6$ 差一倍。]

共形结构有唯一的延拓以填满一个穿刺点。一个有标号的穿刺点定义
$\overline{\cC}_{g,n}/\MCG\to\Mbar_{g,n}$ 的一个截面 $s$。截面满足
$\pi\circ s=id$，微分得 $\dd\pi\circ\dd s=\dd\,id$。在节点处，$\dd\pi$ 在
节点张开的方向上消失——对 $\pi(z,w)=t$，则 $\dd\pi=w\dd z+z\dd w$ 在原点处
为零。因此 $\Mbar$ 上 $\overline{\cC}/\MCG$ 的截面与节点不相交。沿着穿刺截面
$s:\Mbar\to\overline{\cC}/\MCG$，我们考虑切线的族 $(\mathrm{Ker}\,\dd\pi)|_s$，
或其对偶族 $(\mathrm{Ker}\,\dd\pi)^*|_s$。在 Carel Faber 的讲义中，该 Chern
类记作 $K$。

-由穿刺截面 $s$ 把对偶族 $(\mathrm{Ker}\,\dd\pi)^*$ 拉回到 $\Mbar$ 上，所得的
是模空间几何的典范 psi 类 $\psi$。-

% [Tutor Note: p.12 这一句在原文里前后各带一个孤立的连字符"-"（已裁 600 dpi
%  复核，确为排版上的两个短横，不是破折号也不是 $\psi$ 的上下标），看起来是作者
%  原稿里标记"待插入/待改写"的痕迹。本译本按原文照录这两个连字符，未作改动；
%  句读应为：把 $(\mathrm{Ker}\,\dd\pi)^*$ 沿穿刺截面拉回，得到典范 psi 类
%  $\psi$，且 $\psi=c_1\big((\mathrm{Ker}\,\dd\pi)^*|_s\big)$。注意 $K$ 与
%  $\psi$ 在原文里指同一个线丛的不同记法（$K$ 为 Carel Faber 讲义的记法）。]

在这些讲义中，为了强调概念与底层的几何，我们有时会非正式地把一个线丛与其 Chern
类互换着说，非正式地把模空间称为流形、把万有曲线称为纤维丛，并且有时在该用到增广
商的地方仍称开模空间。我们的目标是讨论核心的问题。尽管处理方式是非正式的，有经验的
读者会发现处理是完整的。

上述材料的基本参考文献是 [Bus92] 与 [Wlp10]。

\para{体积结果概览。}

Mirzakhani 证明：WP 体积 $V_{g,n}(L)=V(\Mgb(L))$ 是 $L$ 的多项式，其系数由
$\kappa_1$ 的幂与 $\psi$ 的幂

% [原文 p.13]
的相交数给出。她进一步证明，她用以确定体积多项式的递推满足
Witten--Kontsevich 猜想的定义性关系。以下定理是这两篇论文的直接结果。

\begin{thm}{4} WP 体积多项式。体积多项式由总次数更小的体积多项式递推确定，
[Mir07a, 公式 (5.1) 与定理 8.1]。亏格 $g$、$n$ 条边界、边界长度为
$L=(L_1,\dots,L_n)$ 的双曲曲面的模空间，其体积 $V_{g,n}(L_1,\dots,L_n)$ 是
一个多项式
\[
V_{g,n}(L)=\sum_{\substack{\alpha\\ |\alpha|\le 3g-3+n}}C_\alpha L^{2\alpha},
\]
其中 $\alpha=(\alpha_1,\dots,\alpha_n)$ 为多重指标，且 $C_\alpha>0$ 落在
$\pi^{6g-6+2n-2|\alpha|}\Q$ 中，[Mir07a, 定理 1.1 与 6.1]。其系数是由下式
给出的相交数
\[
C_\alpha=\frac{2^{\delta_{1g}\delta_{1n}}}{2^{|\alpha|}\alpha!\,(3g-3+n-|\alpha|)!}
\int_{\Mgnb}\psi_1^{\alpha_1}\cdots\psi_n^{\alpha_n}\,\omega^{3g-3+n-|\alpha|},
\]
其中 $\psi_j$ 是沿第 $j$ 个穿刺点的余切线的 Chern 类，$\omega$ 是辛形式，
$\alpha!=\prod_{j=1}^n\alpha_j!$，而 $\delta_{**}$ 是 Kronecker 指示 delta，
[Mir07b, 定理 4.4]。\end{thm}

% [Tutor Note: p.13 定理 4 里 $L^{2\alpha}$ 是多重指标的简写
%  $\prod_{i=1}^n L_i^{2\alpha_i}$；因此 $\alpha_i$ 只能是整数，且次数
%  $2|\alpha|=2\sum_i\alpha_i\le 6g-6+2n$ 正好是体积作为边界长度多项式的总次数
%  （与 p.3 的同一陈述一致）。$\delta_{**}$ 是"Kronecker 指示 delta"的简写，
%  实指 $\delta_{1g}\delta_{1n}$：只有 $(g,n)=(1,1)$ 时分子为 $2$，其余情形为
%  $1$。这多出来的因子 $2$ 只在 $(1,1)$ 情形起作用；上面关于 $V_{1,1}$ 的批注
%  给出了它的一次数值印证（$(1,1)$ 时 $C_{(0)}=2\int\omega$）。]

\begin{thm}{5} [Mir07a, 定理 6.3 与 6.4]。体积多项式首项系数的递推关系。对多重
指标 $\alpha$，定义
\[
(\alpha_1,\dots,\alpha_n)_g=C_\alpha\times2^{-\delta_{1g}\delta_{1n}}
\times\prod_{i=1}^n\alpha_i!\times2^{|\alpha|},
\]
则当 $n>0$ 且 $\sum_i\alpha_i=3g-3+n$ 时，有 dilaton 方程
\[
(1,\alpha_1,\dots,\alpha_n)_g=(2g-2+n)(\alpha_1,\dots,\alpha_n)_g,
\]
而当 $n>0$ 且 $\sum_i\alpha_i=3g-2+n$ 时，有 string 方程
\[
(0,\alpha_1,\dots,\alpha_n)_g=\sum_{\alpha_i\neq0}(\alpha_1,\dots,\alpha_i-1,\dots,\alpha_n)_g.
\]
对相交数生成函数
\[
\mathbf{F}(\lambda,t_0,t_1,\dots)=\sum_{g=0}^\infty\lambda^{2g-2}
\sum_{\{d_j\}}\langle\prod_{j=1}^\infty\tau_{d_j}\rangle_g\prod_{r\ge0}t_r^{n_r}/n_r!,
\]
其中 $n_r=\#\{j\mid d_j=r\}$，以及
\[
\langle\tau_{d_1}\cdots\tau_{d_n}\rangle_g=\int_{\Mgnb}\prod_{j=1}^n\psi_j^{d_j},
\]
当 $\sum_{j=1}^n d_j=3g-3+n$ 时成立、而在其他情形该乘积 $\langle\tau_*\rangle$
取为零，则指数 $e^{\mathbf{F}}$ 满足 Virasoro 代数约束，
[Mir07b, 定理 6.1]。\end{thm}

% [Tutor Note: p.13 定理 5 的定义式其实是把定理 4 的 $C_\alpha$ 反解出来的
%  "归一化括号"。代入 $C_\alpha$ 后所有常数全部约掉，得到
%  $(\alpha_1,\dots,\alpha_n)_g=\int_{\Mgn}\psi^{\alpha}\,
%  \omega^{3g-3+n-|\alpha|}/(3g-3+n-|\alpha|)!$。这条恒等式给出三个可当场核对
%  的锚点：(i) 当 $|\alpha|=3g-3+n$ 时 $(3g-3+n-|\alpha|)!=0!=1$、$\omega^0=1$，
%  于是 $(\alpha)_g=\int\psi_1^{\alpha_1}\cdots=\langle\tau_{\alpha_1}\cdots
%  \tau_{\alpha_n}\rangle_g$，与本节末尾 Witten 的括号记号一致；
%  (ii) $|\alpha|=0$ 时 $(0,\dots,0)_g=\int\omega^{3g-3+n}/(3g-3+n)!$，对
%  $(g,n)=(1,1)$ 给 $\pi^2/12$，而定理 4 的常数项
%  $C_{(0)}=2^{1}\cdot\pi^2/12=\pi^2/6$，与 $V_{1,1}(L)=\pi^2/6+L^2/24$ 吻合；
%  (iii) dilaton 方程两边的标记数相差 $1$：左端有 $n+1$ 个指标且总次数
%  $1+\sum_i\alpha_i=3g-2+n=3g-3+(n+1)$，恰好落在"最高次项"的判别式上。]
